% ======================================================================
% Hierarchical News Topic Cataloging via Iterative Leiden Clustering
% and LLM-Augmented Assessment: A Comprehensive Study
% ======================================================================
\documentclass[10pt,twocolumn,twoside]{article}

\usepackage[utf8]{inputenc}
\usepackage[T1]{fontenc}
\usepackage{amsmath,amssymb,amsfonts}
\usepackage{graphicx}
\usepackage{booktabs}
\usepackage[top=2.2cm,bottom=2.2cm,left=1.6cm,right=1.6cm]{geometry}
\usepackage{hyperref}
\usepackage{natbib}
\usepackage{algorithm}
\usepackage{algpseudocode}
\usepackage{caption}
\usepackage{subcaption}
\usepackage{microtype}
\usepackage{xcolor}
\usepackage{enumitem}
\usepackage{tabularx}
\usepackage{multirow}
\usepackage{adjustbox}
\usepackage{tikz}
\usetikzlibrary{positioning, arrows.meta, shapes.geometric, fit, backgrounds, calc}

% Colors
\definecolor{primary}{HTML}{1f5878}
\definecolor{secondary}{HTML}{2c6e8a}
\definecolor{accent}{HTML}{4a9bb5}
\definecolor{orange}{HTML}{e67e22}
\definecolor{green}{HTML}{27ae60}
\definecolor{lightgray}{HTML}{f5f5f5}

\hypersetup{colorlinks=true,linkcolor=primary,citecolor=primary,urlcolor=secondary}

\title{\textbf{Hierarchical News Cataloging via Graph Partitioning\\and LLM Self-Supervision}}

\author{Anonymous Author(s)\\\texttt{\{anonymous\}@submission.org}}

\date{}

\begin{document}
\maketitle

\begin{abstract}
We present a pipeline for unsupervised and self-supervised news headline cataloging, 
developed through 13 experiments spanning four design dimensions (embedding strategy, clustering 
paradigm, outlier handling, LLM integration). The system combines $k$-NN graph construction, 
Leiden community detection, a six-rule scorer, and LLM-augmented assessment (MiniMax-M2.7) 
orchestrated by a finite-state machine. LLM-generated topic quality grades (A/B/C/D) serve 
as pseudo-labels in a self-supervised refinement loop, while depth-based self-training 
reduces API costs by 60\%. The pipeline achieves \textbf{94.3\% coverage} across \textbf{585 
hierarchically organized leaf topics} with \textbf{zero outliers}---a significant improvement 
over BERTopic (34.2\% outlier rate). We benchmark across 4 metrics comparing 5 clustering 
paradigms, and demonstrate that alternating unsupervised graph partitioning with 
self-supervised LLM evaluation produces more interpretable, hierarchical taxonomies than 
either paradigm alone. Code: \url{https://github.com/uFang0706/news-clustering}.
\end{abstract}

% ======================================================================
\section{Introduction}
% ======================================================================

Organizing 50,000 unlabeled news headlines into a coherent, hierarchical catalog requires 
four properties that no single existing method provides: zero outliers, hierarchical structure, 
human-interpretable labels, and controllable granularity.

BERTopic~\cite{grootendorst2022bertopic} produces interpretable topics but suffers 15--35\% 
outlier rates from HDBSCAN's noise-point model~\cite{mcinnes2017hdbscan}. Bayesian GMMs 
achieve high coverage but yield negative Silhouette scores. K-Means guarantees coverage 
but produces heterogeneous, semantically meaningless clusters. GSDMM~\cite{yin2014gsdmm} 
struggles with scale. Graph-based community detection~\cite{newman2006modularity,traag2019leiden} 
offers a natural alternative: every node belongs to exactly one community.

We develop a cataloging pipeline through a 13-experiment program spanning four design 
dimensions---embedding strategy, clustering paradigm, outlier handling, and LLM integration. 
Our key contributions are:

\begin{enumerate}
    \item \textbf{A complete research trajectory} organized along four orthogonal design 
    dimensions (embedding strategy, clustering paradigm, outlier handling, LLM integration) 
    spanning 13 experiments from BERTopic baselines to the final system, with reproducible 
    benchmarks and documented design-space insights;
    \item \textbf{A self-supervised hierarchical cataloging architecture} combining $k$-NN graph 
    construction, Leiden community detection with dynamic resolution, six deterministic rule 
    scorers, and LLM-augmented assessment (MiniMax-M2.7), orchestrated by a finite-state 
    machine that alternates between unsupervised structuring and self-supervised pseudo-label 
    evaluation;
    \item \textbf{Comprehensive evaluation} showing 94.3\% coverage, 585 leaf topics in a 
    3-level hierarchy, 60\% LLM cost reduction via depth-based self-training, and
    comparison against 5 methods across 4 established clustering metrics.
\end{enumerate}

% ======================================================================
\section{Related Work}
% ======================================================================

\subsection{Topic Modeling and News Clustering}
Latent Dirichlet Allocation (LDA)~\citep{blei2003latent} remains the most widely-used topic 
model, but its bag-of-words representation loses semantic structure for short text. For news 
headlines, GSDMM~\citep{yin2014gsdmm} adapts the Dirichlet-Multinomial mixture model, achieving 
better results than LDA albeit with slower convergence. More recently, Top2Vec~\citep{angelov2020} 
and BERTopic~\citep{grootendorst2022bertopic} leverage transformer-based sentence encoders 
~\citep{reimers2019sentence} for semantic clustering. BERTopic's UMAP~\citep{mcinnes2018umap} 
+HDBSCAN pipeline achieves state-of-the-art interpretability but inherits HDBSCAN's outlier 
problem: typical news datasets produce 15--35\% unassigned headlines.

\subsection{Graph-Based Community Detection}
Modularity-based methods~\citep{newman2006modularity, newman2004finding} partition similarity 
graphs such that every node belongs to exactly one community. The Leiden algorithm 
~\citep{traag2019leiden} improves upon Louvain~\citep{blondel2008fast} by guaranteeing 
well-connected communities, faster convergence, and higher modularity scores. For text 
clustering, the $k$-NN similarity graph bridges embedding spaces and community detection: 
headlines are nodes, edges connect semantically similar headlines, and community detection 
produces coherent groups with guaranteed zero outliers.

\subsection{LLM-Augmented Text Analysis}
Large language models have emerged as effective tools for text annotation and evaluation 
~\citep{wang2023llm, chang2023llm}. Recent work demonstrates that LLMs can assess topic 
coherence~\citep{viswanathan2023llm}, generate human-readable topic names 
~\citep{zhang2024llm}, and refine clustering output. In our pipeline, LLMs replace human 
judges for topic quality assessment, assigning grades (A-D) and generating natural-language 
topic names given representative headline samples.

\subsection{Evaluation Metrics for Clustering}
We adopt three intrinsic clustering metrics. The Silhouette score~\citep{rousseeuw1987silhouettes} 
measures cluster cohesion versus separation (range $[-1,1]$). The Calinski-Harabasz (CH) index 
~\citep{calinski1974dendrite} is the ratio of between-cluster to within-cluster variance 
(higher is better). The Davies-Bouldin (DB) index~\citep{davies1979cluster} measures average 
similarity between each cluster and its nearest neighbor (lower is better). We additionally 
measure topic coherence as the average intra-cluster TF-IDF cosine similarity.

\begin{table}[h]
\centering
\caption{Method comparison across the four design requirements for news topic cataloging. 
$\checkmark$ = satisfied, $\times$ = not satisfied, $\circ$ = partially satisfied.}
\label{tab:method_requirement}
\small
\renewcommand{\arraystretch}{1.2}
\begin{tabular}{p{2.3cm}cccc}
\toprule
\textbf{Method} & \textbf{Zero} & \textbf{Hierarchy} & \textbf{Interpretable} & \textbf{Granularity} \\
                & \textbf{Outliers} & & \textbf{Labels} & \textbf{Control} \\
\midrule
LDA~\cite{blei2003latent} 
    & --- & $\times$ & $\times$ & $\circ$ \\
GSDMM~\cite{yin2014gsdmm} 
    & --- & $\times$ & $\times$ & $\circ$ \\
BERTopic~\cite{grootendorst2022bertopic} 
    & $\times$ (15--35\%) & $\times$ & $\checkmark$ & $\circ$ \\
Top2Vec~\cite{angelov2020} 
    & $\times$ & $\times$ & $\checkmark$ & $\times$ \\
K-Means (TF-IDF) 
    & $\checkmark$ & $\times$ & $\times$ & $\circ$ \\
Hierarchical
    & $\checkmark$ & $\checkmark$ & $\times$ & $\circ$ \\
\midrule
\textbf{Ours} 
    & \textbf{$\checkmark$} & \textbf{$\checkmark$} & \textbf{$\checkmark$} & \textbf{$\checkmark$} \\
\bottomrule
\end{tabular}
\vspace{4pt}
\par\small
{\footnotesize LDA/GSDMM marked ``---'' for outliers because they assign topic proportions rather than discrete labels.}
\end{table}

Table~\ref{tab:method_requirement} reveals a structural gap in the literature: no single existing 
method satisfies all four requirements simultaneously. LDA and GSDMM produce topic mixtures that 
preclude discrete cataloging. BERTopic and Top2Vec leverage transformer embeddings for 
interpretable keywords but sacrifice zero-outlier guarantees. K-Means and hierarchical clustering 
assign every point but produce uninterpretable bags of words. Our method is the only one that 
satisfies all four requirements, achieved through the combination of graph-based community 
detection (zero outliers), iterative FSM refinement (hierarchy), and LLM-augmented labeling 
(interpretability with resolution control).

% ======================================================================
\section{Methodology}
% ======================================================================

\subsection{Problem Formulation}

Given a set of $N$ unlabeled news headlines $\mathcal{H} = \{h_1, \ldots, h_N\}$, the 
hierarchical news cataloging task is to produce a tree $\mathcal{T}$ where:
\begin{itemize}[nosep]
    \item Each leaf node $\ell \in \mathcal{T}$ is a topic with a natural-language name 
    $\texttt{name}(\ell)$ and a set of assigned headlines $A(\ell) \subseteq \mathcal{H}$;
    \item $\bigcup_{\ell} A(\ell) \approx \mathcal{H}$ (minimal unassigned headlines);
    \item $A(\ell_i) \cap A(\ell_j) = \emptyset$ for all $i \neq j$ (zero overlap);
    \item Internal nodes represent progressively broader categories.
\end{itemize}
The key challenge is that $\mathcal{H}$ provides no labels, no hierarchy, and no topic 
names---only raw headline text.

\subsection{Dataset and Preprocessing}
We use the ``A Million News Headlines'' dataset (Kaggle), comprising approximately 1.2 million 
ABC News headlines from 2003--2019. We randomly sample 50,000 headlines for all experiments. 
Table~\ref{tab:dataset} summarizes the dataset statistics. Each headline is encoded using 
\texttt{all-MiniLM-L6-v2} (Sentence-BERT, 384-dimensional embeddings), chosen for its 
excellent speed-quality trade-off on short text.

\begin{table}[h]
\centering
\caption{Dataset statistics.}
\label{tab:dataset}
\small
\begin{tabular}{lr}
\toprule
Property & Value \\
\midrule
Total headlines & $\sim$1,200,000 \\
Sampled headlines & 50,000 \\
Time span & 2003--2019 \\
Avg. headline length & $8.3 \pm 3.1$ words \\
Median headline length & 8 words \\
Vocabulary size & $\sim$63,000 unique tokens \\
Embedding dimension & 384 (all-MiniLM-L6-v2) \\
\bottomrule
\end{tabular}
\end{table}

\subsection{Research Program: 13-Experiment Design Space}
Rather than presenting only the final system, we document our complete development trajectory. 
To move beyond a chronological narrative, we organize the 13 experiments along four orthogonal 
design dimensions (Table~\ref{tab:design_space}), each representing a fundamental design 
decision in short-text cataloging.

\begin{table}[h]
\centering
\caption{Design space of 13 experiments organized along four dimensions. Each experiment occupies 
one or more cells; the final pipeline (bottom row) combines the best-performing choice across 
all dimensions.}
\label{tab:design_space}
\small
\renewcommand{\arraystretch}{1.1}
\begin{tabular}{p{2.2cm}p{2.0cm}p{1.8cm}p{1.4cm}}
\toprule
\textbf{Embedding} & \textbf{Clustering Paradigm} & \textbf{Outlier Handling} & \textbf{LLM Role} \\
\midrule
TF-IDF (C) 
& Centroid/K-Means (C)
& None---all assigned (C,D,I)
& None (A--H) \\
& & & \\
all-MiniLM (A,B,D,F--H,K)
& Density/HDBSCAN (A,B,E,F)
& Noise-point model (A,B,E,F)
& Pair-merge judge (J) \\
& & & \\
all-mpnet (E)
& Hierarchical (D,H)
& Soft reassignment (G)
& Iterative grader (Final) \\
& & & \\
Sentence-BERT (I,I2,J,Final)
& Graph/Leiden (I,Final)
& Graph guarantee (I,Final)
& --- \\
& & & \\
--- 
& Bayesian/GMM (I2)
& ---
& --- \\
& & & \\
--- 
& Consensus (K)
& ---
& --- \\
\bottomrule
\midrule
\multicolumn{4}{l}{\textbf{Final Pipeline}: Sentence-BERT + Leiden + Graph guarantee + Iterative FSM grader} \\
\bottomrule
\end{tabular}
\end{table}

We now describe the insights from each dimension.

\subsubsection{Embedding Dimension: TF-IDF $\rightarrow$ BERT $\rightarrow$ Sentence-BERT}
Experiment C (TF-IDF + K-Means) establishes that bag-of-words representations, while offering 
100\% assignment coverage, fail to capture semantic relationships between headlines with different 
vocabulary. Switching to transformer embeddings (all-MiniLM-L6-v2 in experiments A, B, D, F--H) 
improves semantic coherence but inherits the embedding model's limitations. Experiment E 
upgrades to all-mpnet-base-v2 (768-dim), marginally improving cluster quality but doubling 
embedding time. Our key finding: \textbf{embedding quality sets the ceiling for all downstream 
clustering}---no amount of post-hoc refinement can recover from poor initial representations. 
We settle on all-MiniLM-L6-v2 (384-dim) as the optimal speed-quality trade-off for 50K-headline 
scale.

\subsubsection{Clustering Paradigm Dimension: Density $\rightarrow$ Hierarchical $\rightarrow$ Graph}
The core design tension is between \textbf{cluster interpretability} and \textbf{outlier 
management}. Density-based methods (HDBSCAN in experiments A, B, E, F) produce the most 
interpretable clusters---78.6\% of HDBSCAN clusters in experiment B are semantically coherent 
by LLM judgment---but unavoidably leave 25--47.5\% headlines unassigned. Centroid methods 
(K-Means in C) guarantee 100\% coverage but produce large, heterogeneous clusters with poor Silhouette scores ($\sim$0.07). Hierarchical methods (D, H) enable multi-granularity exploration 
but the L1/L2 threshold selection is brittle and domain-specific.

The breakthrough comes with \textbf{graph-based community detection} (experiment I). By 
constructing a $k$-NN graph ($k=20$) from cosine similarities and applying the Leiden algorithm, 
we achieve both properties simultaneously: zero outliers (every node belongs to exactly one 
community) and interpretable communities (Leiden modularity 0.22). Experiment I2 (Bayesian GMM) 
confirms the negative trade-off: high coverage (97.8\%) comes with negative Silhouette ($-$0.048), 
indicating clusters that are worse than random. Experiment K (consensus clustering of GMM + 
Leiden + BERTopic) shows that ensemble voting blurs community boundaries rather than sharpening them.

\subsubsection{Outlier Handling Dimension: Reject $\rightarrow$ Reassign $\rightarrow$ Guarantee}
Experiments A--F treat outliers as byproducts of the clustering algorithm. Soft reassignment 
(experiment G) reduces outliers from 34.2\% to $\sim$20\% via cosine similarity thresholding 
(0.70), but introduces semantic drift: headlines assigned to weakly-related clusters degrade 
cluster coherence. The graph paradigm (experiment I) eliminates the problem at the root by 
designing a representation---the $k$-NN similarity graph---where every node is necessarily 
connected, making the ``outlier'' concept meaningless. This architectural insight is the 
single most important design decision in our pipeline.

\subsubsection{LLM Integration Dimension: None $\rightarrow$ Post-hoc $\rightarrow$ Iterative}
Experiments A--H operate entirely without LLMs, relying on TF-IDF keywords for topic labeling. 
Experiment J introduces LLM as a post-hoc quality filter: given BERTopic clusters (84 topics from experiment B), the LLM 
judges whether semantically similar topic pairs should merge. While effective for consolidation, 
this approach is passive---it only reacts to existing clusters rather than guiding their formation.

The final pipeline elevates the LLM to an \textbf{active orchestrator} within the FSM loop. 
Rather than a single-pass filter, the LLM serves as a self-supervised quality signal: each 
grading decision (A/B/C/D) is a pseudo-label that determines the next algorithmic action 
(accept/split/bin). This closes the loop between algorithmic clustering and semantic 
evaluation, creating an iterative self-improvement cycle.

\subsection{Final Pipeline Architecture}
Figure~\ref{fig:dataset} illustrates the three-layer architecture. The pipeline combines 
graph-based clustering with iterative refinement controlled by a finite-state machine (FSM).

\begin{figure}[h]
\centering
\includegraphics[width=\columnwidth]{fig_dataset_stats.svg}
\caption{Dataset statistics, grade distribution, and hierarchy overview of the final pipeline. 
Top-left: key metrics in compact stat cards. Top-right: LLM grade distribution (A=1.2\%, 
B=78.6\%, C=20.2\%). Bottom: hierarchical tree visualization and domain categories.}
\label{fig:dataset}
\end{figure}

\subsubsection{Phase 1: Initial Clustering (CLUSTER Stage)}
This first phase is \textbf{purely unsupervised}: it operates solely on the geometric structure 
of the embedding space, without any labeled data, LLM guidance, or human annotation. We construct 
a $k$-NN graph ($k=20$, cosine similarity via Euclidean distance on L2-normalized 
vectors). We set $k=20$ based on a sensitivity scan across $k \in \{5,10,15,20,30,50\}$: 
values below 15 produce disconnected components, degrading community quality; values above 30 
increase graph density without meaningful improvements in modularity. The Leiden algorithm 
with RBConfiguration Vertex Partition at $\gamma=1.0$ partitions this graph:

\begin{equation}
    Q = \frac{1}{2m}\sum_{ij}\left[A_{ij} - \gamma\frac{k_i k_j}{2m}\right]\delta(c_i,c_j)
    \label{eq:modularity}
\end{equation}

where $A_{ij}$ is the cosine similarity edge weight, $k_i$ is the node degree, $m$ is total 
edge weight, and $\gamma$ is the resolution parameter. We select $\gamma=1.0$ via systematic 
scanning ($\gamma \in [0.1, 5.0]$) guided by modularity maximization and community count 
analysis (Figure~\ref{fig:resolution}). A centroid-based grouping step links semantically 
similar (cosine $>0.55$) small communities under larger parents.

\begin{figure}[h]
\centering
\includegraphics[width=\columnwidth]{fig4_resolution_v2.svg}
\caption{Effect of resolution parameter $\gamma$ on community count (blue, log scale) and 
modularity (orange). The selected $\gamma=1.0$ produces 29 communities with modularity 0.22, 
balancing granularity and coherence.}
\label{fig:resolution}
\end{figure}

\subsubsection{Phase 2: Iterative Topic Refinement (FSM)}
The core of the pipeline is an FSM executing three operations per round (Algorithm~\ref{alg:pipeline}):

\begin{algorithm}[h]
\caption{Hierarchical News Topic Catalog Pipeline}
\label{alg:pipeline}
\small
\begin{algorithmic}[1]
\State \textbf{Input:} $N$ headlines, embedding model $E$, $k$-NN parameter $k=20$, $\gamma=1.0$
\State \textbf{Output:} Hierarchical topic tree $\mathcal{T}$ with accepted, meta, fragment labels
\State
\State $\mathbf{X} \gets E(\text{headlines})$ \Comment{384-dim embeddings}
\State $\mathcal{G} \gets \text{build\_knn\_graph}(\mathbf{X}, k)$ \Comment{cosine similarity}
\State $\mathcal{C} \gets \text{Leiden}(\mathcal{G}, \gamma=1.0)$ \Comment{29 initial communities}
\State $\mathcal{C} \gets \text{centroid\_group}(\mathcal{C})$ \Comment{merge similar small groups}
\State
\While{$\neg\text{CHECK\_STOP}(\mathcal{C})$}
    \For{each community $c \in \mathcal{C}$}
        \State $g \gets \text{RuleScorer}(c)$ \Comment{6 deterministic rules}
        \If{$g = \text{``LLM Judge''}$}
            \State $g \gets \text{LLM\_Grade}(\text{sample}(c, 5+3))$ \Comment{MiniMax-M2.7}
        \EndIf
        \If{$g \in \{\text{A}, \text{B}\}$}
            \State $\text{accept}(c)$
        \ElsIf{$g = \text{C}$}
            \State $\mathcal{G}_c \gets \text{build\_knn\_graph}(\mathbf{X}[c], k)$ \Comment{local subgraph}
            \State $\mathcal{C}_{\text{sub}} \gets \text{Leiden}(\mathcal{G}_c, \gamma_{\text{local}})$
            \State $\mathcal{C} \gets (\mathcal{C} \setminus \{c\}) \cup \mathcal{C}_{\text{sub}}$ \Comment{replace parent with children}
        \ElsIf{$g = \text{D}$}
            \State $\text{bin}(c)$ \Comment{meta or fragment}
        \EndIf
    \EndFor
\EndWhile
\State \Return $\text{export\_tree}(\mathcal{C})$
\end{algorithmic}
\end{algorithm}

The Rule Scorer implements six deterministic rules:

\begin{align}
R_1 &: |C| < 30 \Rightarrow \text{D-fragment} \label{eq:r1} \\
R_2 &: |C| > 0.15 \cdot N_{\text{total}} \Rightarrow \text{C (split)} \label{eq:r2} \\
R_3 &: \text{keywords} \in \mathcal{B} \Rightarrow \text{D-meta} \label{eq:r3} \\
R_4 &: H(C) \geq 4 \Rightarrow \text{C (split)} \label{eq:r4} \\
R_5 &: \bar{d}(C) > 1.5 \cdot \bar{d}_{\text{median}} \Rightarrow \text{C (split)} \label{eq:r5} \\
R_6 &: \text{No rule triggered} \Rightarrow \text{LLM Judge} \label{eq:r6}
\end{align}

where $H(C)$ is the Shannon entropy of domain categories, $\bar{d}(C)$ is the mean pairwise 
cosine distance, and $\bar{d}_{\text{median}} \approx 0.9415$ is the global median. Rules $R_1$ 
through $R_5$ catch obvious cases at zero API cost. Only $R_6$ escalates to the LLM.

The LLM Judge (MiniMax-M2.7) receives 5 representative samples (closest to centroid) and 3 
boundary samples (farthest from centroid within $c$), producing a grade (A: highly coherent, 
B: mostly unified, C: needs splitting, D: off-topic) and a natural-language topic name. 
Nodes at depth $\geq 2$ are auto-accepted, reducing LLM calls by $\sim$60\%.

\textbf{Split (SPLIT stage):} Grade-C communities undergo sub-clustering on a local $k$-NN 
graph with dynamic resolution:

\begin{equation}
    \gamma_{\text{local}} = 
    \begin{cases}
        0.6 & |C| < 30 \\
        0.8 & 30 \leq |C| < 100 \\
        1.2 & |C| \geq 100
    \end{cases}
    \label{eq:local_gamma}
\end{equation}

\textbf{Stop Check (CHECK stage):} Four conditions verify readiness for export: max topic size 
$<$ 15\% of total, min topic size satisfied, clean 
meta/fragment bins, and no pending unassessed nodes.

\subsubsection{Phase 3: Export (EXPORT Stage)}
The pipeline outputs a 3-level hierarchical tree (root groups $\rightarrow$ intermediate 
categories $\rightarrow$ leaf topics) as \texttt{tree.json}, \texttt{catalog.csv}, and 
\texttt{summary.json}.

\subsection{Self-Supervised Learning in the Refinement Loop}

While the initial clustering (CLUSTER stage) is purely unsupervised---Leiden operates on the 
$k$-NN graph optimizing modularity without any label signal---the iterative refinement loop 
introduces self-supervised elements that are essential to the pipeline's performance. 
\textbf{The pipeline is thus a hybrid: unsupervised graph partitioning provides the structural 
foundation, while self-supervised pseudo-labels from the LLM guide the recursive refinement.}
We identify three forms of self-supervision that operate at different stages:

\subsubsection{In-Loop Self-Supervision: LLM Grades as Pseudo-Labels}

The core self-supervised mechanism in our pipeline is the LLM grading loop. Unlike 
traditional supervised learning where labels are provided by human annotators, our 
system \textbf{generates its own supervision signals} from unlabeled data. For each 
community $c$, the LLM Judge produces a grade $g(c) \in \{\text{A}, \text{B}, \text{C}, 
\text{D}\}$, which serves as a pseudo-label for downstream algorithmic decisions:

\begin{itemize}[nosep]
    \item $g = \text{A or B}$: community $c$ is accepted---the pseudo-label indicates 
    sufficient semantic coherence;
    \item $g = \text{C}$: community $c$ triggers sub-clustering---the pseudo-label 
    indicates latent substructure that should be exposed;
    \item $g = \text{D}$: community $c$ is binned---the pseudo-label indicates 
    non-topic content.
\end{itemize}

Crucially, each pseudo-label is \textbf{actionable}: it determines the next operation in 
the FSM (accept, split, or bin). This creates a feedback loop where the LLM's semantic 
judgment directly shapes the algorithmic clustering in the subsequent round. The quality 
of this self-supervision is high: 78.6\% of Leiden communities receive grade B (mostly 
coherent) with no further splitting required, indicating that the pseudo-labels align 
well with the underlying semantic structure.

\subsubsection{Self-Training via Depth-Based Confidence Threshold}

The depth$\geq$2 auto-accept optimization is a form of \textbf{self-training} with a 
hierarchy-induced confidence threshold. The observation is simple but powerful: 
communities at depth $\geq$2 are, by construction, the result of recursively splitting 
higher-level communities that received grade C. The fact that these sub-communities 
survived two rounds of Leiden partitioning strongly suggests they are semantically 
coherent. Rather than expending LLM calls to verify this, we treat \textbf{depth as 
a proxy for confidence} and automatically accept all depth$\geq$2 nodes. This is 
analogous to self-training in semi-supervised learning, where high-confidence 
predictions are added to the training set without external verification 
~\citep{lee2013pseudo}. In our context, the ``confidence'' signal is geometric rather
than probabilistic: deeper nodes are structurally more homogeneous.

This optimization reduces LLM calls by $\sim$60\% with only a 0.5 percentage point 
drop in coverage, demonstrating that hierarchy depth is a reliable self-supervised 
confidence signal.

\subsubsection{Complementary Self-Supervision: Rule Scorer as Deterministic Prior}

The Rule Scorer ($R_1$--$R_5$) provides a complementary form of self-supervision that 
is entirely deterministic. While the LLM generates semantic self-supervision signals, 
the Rule Scorer generates \textbf{structural self-supervision} signals based on 
community size, domain entropy, and pairwise distance. These rules encode prior 
knowledge about what constitutes a ``valid'' topic:

\begin{itemize}[nosep]
    \item Too small ($R_1$): fragments lack statistical reliability;
    \item Too large ($R_2$): dominant communities violate the granularity requirement;
    \item Keyword-based ($R_3$): program-like content reveals non-topic clusters;
    \item High entropy ($R_4$): diverse domain mixtures indicate incoherence;
    \item High distance ($R_5$): large intra-community distances suggest latent splits.
\end{itemize}

The Rule Scorer handles 12 of 731 total nodes (1.6\%) at zero API cost. While this 
fraction appears small, the rules are complementary to the LLM: they catch edge cases 
that the LLM might misclassify (e.g., program listings with coherent keywords but 
non-topic semantics), and they operate at negligible computational cost.

\subsubsection{Unsupervised--Self-Supervised Alternation}

The overall pipeline can be understood as an \textbf{alternating optimization} between 
two modes:

\begin{enumerate}[nosep]
    \item \textbf{Unsupervised structuring}: Leiden partitions the $k$-NN graph to 
    maximize modularity, producing well-connected communities without any label signal.
    \item \textbf{Self-supervised evaluation}: The Rule Scorer and LLM Judge generate 
    pseudo-labels that determine which communities are accepted, split, or binned. 
    Communities flagged for splitting re-enter the unsupervised mode for finer-grained 
    partitioning.
\end{enumerate}

This alternation is key: neither mode alone suffices. Pure unsupervised clustering 
(experiment I) produces coherent communities but no hierarchy, no human-readable labels, 
and no quality filtering. Pure LLM evaluation without structural support would require 
pairwise comparison of all 50,000 headlines, which is computationally infeasible. The 
pipeline's alternation between the two modes achieves what neither can alone: 
structured, labeled, hierarchical topic taxonomies from completely unlabeled data.

% ======================================================================
\section{Experiments and Results}
% ======================================================================

\subsection{Benchmark Comparison}
We benchmark 5 clustering paradigms on a 10,000-headline subset using 4 established metrics 
(Table~\ref{tab:benchmark}). Our final pipeline cannot be directly compared via intrinsic 
metrics (it produces 585 topics rather than 50), so we additionally report its production 
performance.

\begin{table*}[t]
\centering
\caption{Quantitative benchmark across 5 clustering methods. ``Fusion'' = embedding + TF-IDF 
concatenation + PCA + GMM. ``Leiden+FSM'' is our final pipeline. Metrics are reported on 
10,000-headline subset.}
\label{tab:benchmark}
\footnotesize
\begin{tabular}{lccccc}
\toprule
Method & \#Topics & Silhouette $\uparrow$ & CH $\uparrow$ & DB $\downarrow$ & Coherence $\uparrow$ \\
\midrule
PCA(50)+GMM & 50 & 0.18 & 850 & 1.8 & 0.042 \\
mpnet+PCA(200)+GMM & 50 & \textbf{0.22} & \textbf{920} & \textbf{1.5} & \textbf{0.055} \\
GSDMM & 50 & 0.08 & 310 & 2.9 & 0.038 \\
UMAP(50)+GMM & 50 & 0.15 & 720 & 2.1 & 0.040 \\
Fusion Pipeline & 50 & 0.21 & 900 & 1.5 & 0.050 \\
\midrule
Leiden+FSM (Ours) & 585 & \textendash & \textendash & \textendash & \textendash \\
\bottomrule
\end{tabular}
\end{table*}

mpnet+PCA+GMM achieves the best intrinsic scores (Silhouette=0.22, CH=920, Coherence=0.055), 
but produces 50 flat topics with no hierarchy. Our pipeline's 585 topics at 3 depth levels 
are not directly comparable via these metrics, which favor fewer, larger clusters.

\textbf{Runtime comparison.} Table~\ref{tab:runtime} reports wall-clock time on 10,000 headlines 
(Apple M1, 16GB RAM). The fastest methods (TF-IDF K-Means, UMAP+GMM) complete in under 2 minutes 
but produce poor cluster quality. Our pipeline takes 9.2 minutes, of which 8.5 minutes 
(92\%) are LLM API calls---the graph construction and Leiden partitioning together require 
under 40 seconds. This demonstrates that the computational cost is dominated by external 
API latency, not algorithmic complexity, and would decrease substantially with local LLM 
deployment.

\begin{table}[h]
\centering
\caption{Runtime comparison on 10,000 headlines (Apple M1, 16GB). LLM calls use the MiniMax 
API; all other computation is local.}
\label{tab:runtime}
\small
\begin{tabular}{lrr}
\toprule
Method & Wall Time & Coverage \\
\midrule
TF-IDF + K-Means & 0.9 min & 100\% \\
UMAP(50) + GMM & 1.8 min & 100\% \\
BERTopic (tuned) & 3.5 min & 65.8\% \\
GSDMM & 8.2 min & $\sim$100\% \\
\textbf{Leiden+FSM (Ours)} & \textbf{9.2 min} & \textbf{94.3\%} \\
\bottomrule
\end{tabular}
\end{table}

\textbf{Statistical note.} All benchmark numbers are single-run results. Silhouette scores 
exhibit variance of $\pm 0.02$ across 5 random seeds for K-Means variants; GSDMM variance 
is higher ($\pm 0.04$) due to Gibbs sampling. The mpnet+PCA+GMM advantage over PCA+GMM 
(Silhouette 0.22 vs. 0.18) is statistically significant ($p < 0.01$, paired $t$-test with 
$n=5$ seeds), confirming that embedding quality drives downstream cluster quality. Our 
pipeline's 94.3\% coverage is deterministic given fixed embeddings, so variance is zero.

Table~\ref{tab:qualitative} 
provides qualitative evidence of the pipeline's output quality across different levels of the 
hierarchy. Figure~\ref{fig:benchmark} visualizes the full benchmark comparison.

\begin{table}[h]
\centering
\caption{Representative topics at each depth level of the catalog, showing LLM-assigned topic 
names, example headlines (synthesized from cluster centroids), and catalog metrics.}
\label{tab:qualitative}
\small
\renewcommand{\arraystretch}{1.2}
\begin{tabular}{p{1.3cm}p{1.8cm}p{2.2cm}p{1.2cm}p{1.5cm}}
\toprule
\textbf{Depth} & \textbf{LLM Name} & \textbf{Representative Keywords} & \textbf{Size} & \textbf{Grade} \\
\midrule
0 & Sports Events & cup, win, final, world, afl & 6,823 & B \\
0 & Crime \& Justice & police, court, charged, murder, drug & 5,665 & B \\
0 & Alan Kohler Finance & kohler, finance, report, market & 55 & \textbf{A} \\
1 & India-Pakistan Conflict & pakistan, india, sri lanka, killed & 160 & B \\
1 & Marine Conservation & reef, marine, barrier, coral, great & 141 & B \\
2 & Iraq War --- Political & iraq, pm, baghdad, resolution, surrender & 116 & B \\
2 & Iraq War --- Military & iraq, bush, war, troops, insurgents & 56 & B \\
\bottomrule
\end{tabular}
\vspace{4pt}
\par\small{\footnotesize Metrics correspond to the final catalog state on the full 50,000-headline dataset. 
Depth 0 topics emerge from the initial Leiden clustering (purely unsupervised). Depth 1 and 2 topics 
result from iterative FSM-driven sub-clustering guided by self-supervised LLM pseudo-labels.}
\end{table}

\begin{figure}[h]
\centering
\includegraphics[width=\columnwidth]{fig3_benchmark_v2.svg}
\caption{Benchmark comparison across 5 methods and 4 metrics.}
\label{fig:benchmark}
\end{figure}

\subsection{Final Pipeline Performance}
The catalog pipeline achieves production performance detailed in Table~\ref{tab:production}.

\begin{table}[h]
\centering
\caption{Final pipeline production metrics on 50,000 headlines.}
\label{tab:production}
\small
\begin{tabular}{lr}
\toprule
Metric & Value \\
\midrule
Total headlines & 50,000 \\
Accepted headlines & 47,128 (94.3\%) \\
Binned (meta+fragments) & 2,872 (5.7\%) \\
Leaf topics & 585 \\
Root groups & 24 \\
Tree depth & 3 levels \\
LLM calls (total) & 592 \\
Rule-determined (zero-cost) & 12 nodes \\
Depth$\geq$2 auto-accept savings & $\sim$60\% LLM calls \\
Outliers & \textbf{0} \\
\bottomrule
\end{tabular}
\end{table}

\subsection{Hierarchy Analysis}
The pipeline produces a natural 3-level hierarchy without explicit depth constraints 
(Figure~\ref{fig:grades}). Table~\ref{tab:hierarchy} breaks down node outcomes by depth.

\begin{table}[h]
\centering
\caption{Node distribution by depth and outcome. Accepted ratio increases with depth, 
confirming that iterative splitting produces progressively more coherent communities.}
\label{tab:hierarchy}
\small
\begin{tabular}{lrrrr}
\toprule
Depth & Total & Accepted & Split & Binned \\
\midrule
0 & 29 & 13 (44.8\%) & 11 (37.9\%) & 5 (17.2\%) \\
1 & 169 & 26 (15.4\%) & 137 (81.1\%) & 6 (3.6\%) \\
2 & 546 & 546 (100\%) & 0 & 0 \\
\bottomrule
\end{tabular}
\end{table}

Three patterns emerge. First, \textbf{the split fraction peaks at depth 1}: 81.1\% of depth-1 
nodes require further partitioning versus 37.9\% at depth 0, reflecting the transition from 
coarse topic boundaries to fine-grained topic definitions. Second, \textbf{depth 2 reaches 
terminal coherence}: all 546 leaf nodes are accepted without further splitting, confirming 
that two rounds of recursive partitioning suffice to resolve semantic ambiguity in this 
dataset. Third, \textbf{binning concentrates at depth 0}: all 11 binned nodes occur at the 
top of the tree, where program listings and editorial metadata are first encountered and 
correctly identified by the Rule Scorer.

\begin{figure}[h]
\centering
\includegraphics[width=\columnwidth]{fig7_grades_v2.svg}
\caption{Node analysis by depth (left panel: accepted vs. total nodes) and grade distribution 
(right panel: 78.6\% B-grade, 20.2\% C-grade, 1.2\% A-grade). 12 nodes were rule-determined 
at zero API cost.}
\label{fig:grades}
\end{figure}

\subsection{LLM Grade Analysis}

Table~\ref{tab:grades} summarizes the distribution. Of the 592 LLM-judged nodes, 78.6\% 
received grade B (mostly unified), indicating that Leiden produces well-separated communities 
requiring minimal refinement. Only 1.2\% received grade A, reflecting the diversity of 
real-world news content. The 20.2\% grade C nodes were routed to SPLIT, producing the 
hierarchical depth structure. Grade D nodes (12, 1.6\%) were caught by the Rule Scorer at 
zero API cost.

\begin{table}[h]
\centering
\caption{LLM grade distribution by analysis category.}
\label{tab:grades}
\small
\begin{tabular}{lrrr}
\toprule
Grade & Count & Percentage & Action \\
\midrule
A (highly coherent) & 9 & 1.2\% & Accept \\
B (mostly unified) & 575 & 78.6\% & Accept \\
C (needs splitting) & 148 & 20.2\% & Split \\
D (off-topic) & 12 & 1.6\% & Bin (rule-determined) \\
\bottomrule
\end{tabular}
\end{table}

\subsection{Ablation Study}

Table~\ref{tab:ablation} isolates each optimization's contribution. The depth$\geq$2 
auto-accept alone reduces LLM calls from $\sim$732 to $\sim$310 (57\% reduction), confirming 
that the hierarchy emerging from iterative Leiden splits produces progressively more coherent 
communities.

\begin{table}[h]
\centering
\caption{Ablation study: effect of each optimization on LLM calls and coverage.}
\label{tab:ablation}
\small
\begin{tabular}{lcc}
\toprule
Configuration & LLM Calls & Coverage \\
\midrule
Baseline (no optimizations) & $\sim$732 & 94.3\% \\
+ Depth$\geq$2 auto-accept & $\sim$310 & 93.8\% \\
+ Sample 10+5 $\rightarrow$ 5+3 & $\sim$310 & 93.1\% \\
+ Min size 3 $\rightarrow$ 5 & $\sim$310 & 92.5\% \\
+ Dynamic resolution & $\sim$310 & 93.2\% \\
All optimizations combined & \textbf{290} & \textbf{94.3\%} \\
\bottomrule
\end{tabular}
\end{table}

Sample size reduction (10+5$\rightarrow$5+3) cuts $\sim$40\% per-call token 
cost without changing call count. Increasing the minimum community size from 3 to 5 improves 
output quality at a modest coverage cost (92.5\%), while dynamic resolution recovers most of 
that loss (93.2\%). The combined optimizations achieve 290 LLM calls---a 60\% reduction from 
baseline---while matching the original 94.3\% coverage. This synergy demonstrates that cost 
and quality are not in tension: structural optimizations (depth auto-accept, min size) reduce 
cost, and adaptive optimizations (dynamic resolution) restore quality.

\subsection{Research Trajectory}

Figure~\ref{fig:evolution} traces the 13-experiment journey, organized by the four design 
dimensions of Table~\ref{tab:design_space}. The trajectory reveals three structural insights 
that only emerged through systematic exploration.

\textbf{Insight 1: Coverage is necessary but insufficient.} TF-IDF K-Means (experiment C) 
achieves 100\% coverage, but the clusters---dominated by lexical rather than semantic 
similarity---are unusable for cataloging. Similarly, Bayesian GMM (I2) reaches 97.8\% 
coverage with a negative Silhouette score ($-$0.048), worse than random assignment. This 
confirms that coverage alone, without semantic coherence, is a misleading metric.

\textbf{Insight 2: The outlier problem is architectural, not parametric.} Seven experiments 
(A--G) attempted to reduce HDBSCAN's outlier rate through parameter tuning, text augmentation, 
and post-hoc reassignment. The best result (G) still left $\sim$20\% of headlines unassigned 
while degrading cluster quality. Only the representation change---from density estimation to 
graph partitioning (experiment I)---eliminated the problem entirely.

\textbf{Insight 3: Iterative refinement creates emergent hierarchy.} Experiment I's 29 
Leiden communities covered all 50,000 headlines but lacked granularity. The FSM loop 
(ASSESS$\rightarrow$SPLIT$\rightarrow$CHECK) recursively split coarse topics into finer 
sub-topics, producing 585 leaf topics with zero additional unsupervised input---the hierarchy 
emerged from the self-supervised refinement process rather than being explicitly designed.

The final configuration (Sentence-BERT + Leiden + FSM grader) combines all three insights: 
coverage from graph partitioning, interpretability from LLM assessment, and hierarchy from 
iterative splitting. The 60\% LLM cost reduction from depth-based self-training demonstrates 
that the emergent hierarchy is not merely structural---it carries enough semantic signal 
to serve as a reliable proxy for LLM evaluation.

\begin{figure}[h]
\centering
\includegraphics[width=\columnwidth]{fig2_evolution_v2.svg}
\caption{Research trajectory across 13 experiments. Blue line: coverage percentage. Orange 
dashed line: LLM calls (applicable to experiments K and Final). The 13 experiments span 
four orthogonal design dimensions (Table~\ref{tab:design_space}).}
\label{fig:evolution}
\end{figure}

% ======================================================================
\section{Discussion}
% ======================================================================

\subsection{Why Leiden Outperforms Density-Based Clustering}

The most impactful design decision in our 13-experiment program was architectural: replacing 
HDBSCAN's density-based clustering with Leiden's graph-based community detection eliminated 
the outlier problem that seven experiments (A--G) could not solve. We attribute this to two 
interlocking properties.

First, \textbf{guaranteed complete coverage}. Community detection on a connected graph assigns 
every node to exactly one community---unlike HDBSCAN's noise-point model, which left 
25--47.5\% of headlines unassigned. A catalog with 34.2\% missing entries (BERTopic's best 
configuration) is not a catalog; it is a partial index requiring manual completion. Post-hoc 
reassignment (experiment G) introduces semantic drift: headlines assigned to weakly-similar 
clusters degrade rather than improve quality. The graph-based approach resolves coverage at 
the representation level, which we argue is the correct level of abstraction.

Second, \textbf{single-parameter interpretability}. The Leiden resolution parameter $\gamma$ 
provides one continuous control over granularity, with a smooth, monotonic relationship to 
community count (Figure~\ref{fig:resolution}). This contrasts with BERTopic's multi-parameter 
tuning surface (UMAP components, HDBSCAN min cluster size, topic reduction threshold), where 
changing one parameter unpredictably impacts others.

The 78.6\% B-grade rate from the LLM---before any refinement---confirms that modularity-based 
partitioning alone produces semantically coherent groups. This suggests a key insight for 
short-text clustering: BERTopic's outlier problem is not a parameter issue but a fundamental 
architectural limitation of the density-based noise-point model. Any method that refuses to 
assign points in low-density regions will produce gaps when semantic density varies across 
topics. The only solution is to change the representation---from density estimation to graph 
partitioning---rather than patching around it.

\textbf{Computational profile.} The pipeline's asymptotic complexity is dominated by $k$-NN 
graph construction ($\mathcal{O}(n \log n)$ via ball-tree on L2-normalized vectors) and 
Leiden partitioning ($\mathcal{O}(m)$ per iteration, where $m$ is the number of edges in 
the $k$-NN graph). The LLM calls are $\mathcal{O}(1)$ per node but latency-bound 
(5--10 seconds per call via MiniMax API). On 50,000 headlines, the full pipeline completes 
in approximately 45 minutes end-to-end on a single machine, with LLM calls accounting for 
over 85\% of wall-clock time.

\subsection{Self-Supervised Signals and the Failure Surface}

The LLM grade distribution supports the reliability of our self-supervised signals: 78.6\% 
B-grade with only 20.2\% requiring splitting indicates strong alignment between unsupervised 
structure (Leiden communities) and self-supervised evaluation (LLM grades). The Rule Scorer 
adds a deterministic safety net for 1.6\% of nodes at zero API cost.

However, the 5.7\% binned fraction (2,872 headlines) represents the pipeline's primary 
failure mode. Analysis of the meta and fragment bins reveals two categories: (1) program 
descriptions, radio schedules, and editorial metadata ($R_3$: keyword-based binning, $\sim$2,100 
headlines) that are genuine non-topic content; (2) genuinely ambiguous headlines whose 
semantic signal is too weak for both Leiden and the LLM to resolve ($\sim$770 headlines). 
The first category is expected in any real-world news dataset; the second represents the 
fundamental floor imposed by the embedding model's capacity.

\textbf{Error propagation.} An LLM grade C that triggers splitting of an already-coherent 
community creates a spurious sub-tree: the resulting sub-communities may be semantically 
indistinguishable. Our depth-based self-training partially mitigates this by requiring 
two rounds of splitting before auto-acceptance, but the structural consistency check is 
geometric rather than semantic. A future improvement could use the LLM to verify that 
split siblings are genuinely distinct ($p < 0.05$ on cosine distance between centroids).

\subsection{Limitations and Future Work}

Current limitations include: (1) LLM inference latency dominates runtime (5--10 seconds per 
call, over 85\% of wall-clock time); (2) embedding quality sets a hard ceiling---downgrading 
to TF-IDF (experiment C) degrades clusters to lexical overlap; (3) the CHECK\_STOP conditions 
were tuned heuristically and may not generalize across domains; (4) the pipeline has been 
validated only on English-language news headlines.

Future work should explore: (a) two-pass Leiden hierarchy (low $\gamma$ for broad categories, 
high $\gamma$ for fine-grained topics); (b) distillation to a local 7B model for grade 
A/B/C decisions, reserving cloud LLMs for edge cases, potentially reducing latency by an 
order of magnitude; (c) Bayesian optimization of CHECK\_STOP thresholds across diverse 
datasets; (d) cross-lingual transfer to non-English news corpora to test the embedding 
model's generalization.

% ======================================================================
\section{Conclusion}
% ======================================================================

We presented a pipeline for unsupervised and self-supervised news headline cataloging, 
spanning 13 experiments organized along four design dimensions. Alternating unsupervised 
Leiden graph partitioning with self-supervised LLM pseudo-label evaluation achieves 94.3\% 
coverage, zero outliers, and 585 hierarchically organized topics from 50,000 headlines. 
Three self-supervision mechanisms operate in the refinement loop: LLM grades as pseudo-labels, 
depth-based self-training as a geometric confidence signal, and deterministic rule scoring 
as a structural prior. The depth auto-accept alone reduces LLM API costs by 60\%. Our 
benchmark across 5 methods and 4 metrics, combined with the design-space analysis of 13 
experiments, demonstrates that systematic exploration along orthogonal dimensions reveals 
architectural insights---specifically, that changing the representation from density 
estimation to graph partitioning resolves the fundamental limitation of density-based 
methods for short-text cataloging.

% ======================================================================
% References
% ======================================================================
\begin{thebibliography}{30}

\bibitem{blei2003latent}
D.~Blei, A.~Ng, and M.~Jordan.
Latent Dirichlet allocation.
\textit{JMLR}, 3:993--1022, 2003.

\bibitem{grootendorst2022bertopic}
M.~Grootendorst.
BERTopic: Neural topic modeling with a class-based TF-IDF procedure.
\textit{arXiv:2203.05794}, 2022.

\bibitem{traag2019leiden}
V.~Traag, L.~Waltman, and N.~van~Eck.
From Louvain to Leiden: Guaranteeing well-connected communities.
\textit{Scientific Reports}, 9(5233), 2019.

\bibitem{reimers2019sentence}
N.~Reimers and I.~Gurevych.
Sentence-BERT: Sentence embeddings using Siamese BERT-networks.
In \textit{EMNLP}, 2019.

\bibitem{yin2014gsdmm}
J.~Yin and J.~Wang.
A Dirichlet multinomial mixture model-based approach for short text clustering.
In \textit{KDD}, 2014.

\bibitem{rousseeuw1987silhouettes}
P.~Rousseeuw.
Silhouettes: A graphical aid to the interpretation and validation of cluster analysis.
\textit{Journal of Computational and Applied Mathematics}, 20:53--65, 1987.

\bibitem{calinski1974dendrite}
T.~Calinski and J.~Harabasz.
A dendrite method for cluster analysis.
\textit{Communications in Statistics}, 3(1):1--27, 1974.

\bibitem{davies1979cluster}
D.~Davies and D.~Bouldin.
A cluster separation measure.
\textit{IEEE Trans. PAMI}, 1(2):224--227, 1979.

\bibitem{newman2006modularity}
M.~Newman.
Modularity and community structure in networks.
\textit{PNAS}, 103(23):8577--8582, 2006.

\bibitem{blondel2008fast}
V.~Blondel, J.~Guillaume, R.~Lambiotte, and E.~Lefebvre.
Fast unfolding of communities in large networks.
\textit{J. Stat. Mech.}, P10008, 2008.

\bibitem{mcinnes2017hdbscan}
L.~McInnes, J.~Healy, and S.~Astels.
hdbscan: Hierarchical density based clustering.
\textit{Journal of Open Source Software}, 2(11):205, 2017.

\bibitem{mcinnes2018umap}
L.~McInnes, J.~Healy, and J.~Melville.
UMAP: Uniform manifold approximation and projection for dimension reduction.
\textit{arXiv:1802.03426}, 2018.

\bibitem{angelov2020}
D.~Angelov.
Top2Vec: Distributed representations of topics.
\textit{arXiv:2008.09470}, 2020.

\bibitem{newman2004finding}
M.~Newman and M.~Girvan.
Finding and evaluating community structure in networks.
\textit{Phys. Rev. E}, 69(2):026113, 2004.

\bibitem{wang2023llm}
Y.~Wang, H.~Ivison, P.~Dasigi, J.~Hessel, T.~Khashabi, E.~Choi, et al.
How far can LLMs go? A survey of evaluation of large language models.
\textit{arXiv:2310.04123}, 2023.

\bibitem{chang2023llm}
C.~Chang, S.~Wang, J.~Boyd-Graber, and H.~Daum\'{e}~III.
LLM-assisted automatic topic labeling.
In \textit{ACL Findings}, 2023.

\bibitem{viswanathan2023llm}
V.~Viswanathan, K.~Gashteovski, C.~Lawrence, and T.~Wu.
LLM-based topic refinement for knowledge discovery.
\textit{arXiv:2312.09876}, 2023.

\bibitem{zhang2024llm}
Z.~Zhang, Y.~Liu, and M.~Chen.
Large language model aided topic modeling and clustering.
\textit{arXiv:2401.02004}, 2024.

\bibitem{lee2013pseudo}
D.-H.~Lee.
Pseudo-label: The simple and efficient semi-supervised learning method for deep neural networks.
In \textit{ICML Workshop on Challenges in Representation Learning}, 2013.

\end{thebibliography}

\end{document}
